\section{Cross-Regime Analysis}
\label{sec:cross}

\subsection{Overall direction}
Matched-six means are positive in all three regimes (\Cref{tab:matched}).
UNSEEN adds a uniformly positive replicate pattern (21/21) on a short technical source.
LOW shows that large positive $\Delta$ can appear under intentionally low thematic overlap.
HIGH is the only regime with negative configuration mean $\Delta$ (both Gemini conditions).
Taken together, the regimes support a tempered summary: positive associations are common across regimes, but HIGH shows that common is not universal at the configuration level.

\subsection{Configuration-specific patterns}
We characterize heterogeneity without ranking:

\begin{itemize}[leftmargin=*]
  \item \textbf{Gemini ON/OFF:} negative or near-zero mean $\Delta$ on HIGH with mixed replicates; strongly positive on LOW; moderately positive on UNSEEN. This is a clear HIGH$\rightarrow$LOW direction reversal.
  \item \textbf{Claude Sonnet 5:} positive mean $\Delta$ in all three regimes; HIGH/LOW large; UNSEEN smaller ($+6.42$\,pp).
  \item \textbf{DeepSeek Expert ON:} persistently positive with attenuation HIGH$\rightarrow$LOW$\rightarrow$UNSEEN ($+9.25\rightarrow +6.19\rightarrow +3.83$\,pp).
  \item \textbf{GPT-5.6 Luna:} persistently positive; HIGH larger than LOW/UNSEEN, which are similar ($+5.31$ / $+5.52$\,pp).
  \item \textbf{Grok 4.5 Fast:} large positive HIGH; much smaller LOW mean with one negative replicate; positive UNSEEN.
  \item \textbf{Qwen 3.8 Max Fast:} positive HIGH and UNSEEN means; no valid LOW paired effect.
\end{itemize}

These patterns differ in kind, not merely in rank order.
Reversal (Gemini), attenuation (DeepSeek), magnitude instability without mean-sign flip (Grok), and missing cells (Qwen LOW) are distinct descriptive phenomena.

\subsection{Overlap is not a sufficient explanation}
If larger overlap were necessary for positive priming $\Delta$, LOW and UNSEEN should not show widespread positives.
They do.
Gemini's largest positives appear on LOW rather than HIGH.
Overlap, length, and genre remain confounded, so these observations reject a simple ``overlap causes $\Delta$'' slogan without identifying an alternative single cause.
A responsible scientific statement is therefore comparative and negative with respect to that slogan: high overlap is not required for positive coverage $\Delta$ in this design.

\subsection{Length and genre as disclosed confounds}
UNSEEN's smaller matched-six mean ($+5.78$\,pp) should not be narrated as ``weaker priming because the domain is unseen'' without acknowledging that UNSEEN is also shorter and expository.
Likewise, HIGH's mixture of large positives and Gemini negatives cannot be attributed solely to overlap.
The cross-regime section therefore prioritizes pattern description over monocausal storytelling.

\subsection{RQ4}
Configuration heterogeneity is first-order: sign reversals, attenuation, amplification, and invalid cells coexist in the same experimental frame.
Collapsing the study into a single universal priming effect would misrepresent the evidence.

\subsection{Practical reading for manuscript users}
For readers who need a one-paragraph takeaway without collapsing heterogeneity, the safest summary is:
matched-six mean $\Delta$ is positive in HIGH, LOW, and UNSEEN; UNSEEN replicate $\Delta$s are uniformly positive; HIGH uniquely contains negative Gemini means; LOW uniquely lacks a valid Qwen paired estimate; qualitative audits show that HIGH positives often involve opening reuse while LOW/UNSEEN positives often do not.
Any shorter slogan loses at least one of those facts.
